\documentclass[a4,11pt]{aleph-notas}

% -- Paquetes adicionales
\usepackage{enumitem}
\usepackage{aleph-comandos}
\usepackage{esint}

% -- Datos
\institucion{Facultad de Ciencias Exactas, Naturales y Ambientales}
\carrera{Catálogo STEM}
\asignatura{Matemática III}
\tema{Resumen 17: Teoremas integrales}
\autor[A. Merino]{Andrés Merino}
\fecha{Periodo 2026-1}

\logouno[0.16\textwidth]{Logos/logoPUCE_04_ac}
\definecolor{colortext}{HTML}{1957d1}
\definecolor{colordef}{HTML}{1957d1}
\fuente{montserrat}


\begin{document}

\encabezado

%%%%%%%%%%%%%%%%%%%%%%%%%%%%%%%%%%%%%%
\section{Teorema de Green}
%%%%%%%%%%%%%%%%%%%%%%%%%%%%%%%%%%%%%%

\begin{teo}[Teorema de Green]
    Sea $\Omega\subseteq \R^2$ un conjunto cerrado, acotado y simplemente conexo tal que $\partial \Omega$ es una curva suave parametrizada por una trayectoria que recorre $\partial \Omega$ de manera antihoraria. Se tiene que
    \[
        \iint_\Omega \rot(f) \,dA = \oint_{\partial \Omega} f\cdot ds.
    \]
\end{teo}

\begin{advertencia}
    Si tenemos que
    \[
        f(x,y)=(M(x,y),N(x,y)),
    \]
    para $(x,y)\in\R^2$, entonces el Teorema de Green se enuncia por
    \[
        \iint_\Omega \parcial{N}{x} - \parcial{M}{y} \,dxdy = \oint_{\partial \Omega} M\, dx+ N\, dy.
    \]
\end{advertencia}

%%%%%%%%%%%%%%%%%%%%%%%%%%%%%%%%%%%%%%
\section{Teorema de Gauss en el plano}
%%%%%%%%%%%%%%%%%%%%%%%%%%%%%%%%%%%%%%

\begin{teo}[Teorema de Gauss en el plano]
    Sea $\Omega\subseteq \R^2$ un conjunto cerrado, acotado y simplemente conexo tal que $\partial \Omega$ es una curva suave parametrizada por una trayectoria $\func{\alpha}{I}{\R^2}$ que recorre $\partial \Omega$ de manera antihoraria. Se define
    \[
        n(t) = \frac{(\alpha_2'(t),- \alpha_1'(t))}{\|\alpha'(t)\|}
    \]
    para $t\in I$. Se tiene que
    \[
        \iint_\Omega \dive(f) \,dA = \oint_{\partial \Omega} f\cdot n\, d\alpha.
    \]
\end{teo}

%%%%%%%%%%%%%%%%%%%%%%%%%%%%%%%%%%%%%%
\section{Teorema de Stokes}
%%%%%%%%%%%%%%%%%%%%%%%%%%%%%%%%%%%%%%

\begin{defi}[Superficie orientable]
    Se dice que una superficie $\Sigma$ es orientable si es posible definir un solo vector normal unitario en cada punto de $\Sigma$ de modo que este sea continuo. Una vez definido el vector normal se dice que $\Sigma$ es una superficie orientada.
\end{defi}

\begin{defi}
    Dada una superficie orientada $\Sigma$ y una curva $C$ sobre $\Sigma$ se dice que la parametrización de $C$ tiene orientación consistente con la parametrización de $\Sigma$ si el vector normal de $\Sigma$ sigue la regla de la mano derecha con respecto a $C$.
\end{defi}

\begin{teo}[Teorema de Stokes]
    Sea $\Sigma$ una superficie orientada, acotada y suave en $\R^3$ tal que el borde de $\Sigma$, denotado por $\partial \Sigma$, es una curva suave, simple, cerrada y parametrizada con orientación consistente con la parametrización de $\Sigma$. Sea $\func{f}{\Sigma}{\R^3}$ un campo vectorial diferenciable, se tiene que
    \[
        \iint_\Sigma \rot(f) \cdot dS = \oint_{\partial \Sigma} f\cdot ds.
    \]
\end{teo}

%%%%%%%%%%%%%%%%%%%%%%%%%%%%%%%%%%%%%%
\section{Teorema de Gauss}
%%%%%%%%%%%%%%%%%%%%%%%%%%%%%%%%%%%%%%

\begin{teo}[Teorema de Gauss]
    Sea $\Omega$ una región acotada de $\R^3$ tal que $\partial \Omega$ es una superficie suave, cerrada y parametrizada con orientación hacia fuera de $\Omega$. Sea $\func{F}{\Omega}{\R^3}$ un campo vectorial diferenciable, se tiene que
    \[
        \iiint_\Omega \dive(F) \, dV = \varoiint_{\partial \Omega} F\cdot dS.
    \]
\end{teo}

\begin{prop}
    Sea $\Omega$ una región acotada de $\R^3$ tal que $\partial \Omega$ es una superficie suave, cerrada y parametrizada con orientación hacia fuera de $\Omega$. Sean $\func{F}{\Omega}{\R^3}$ un campo vectorial y $\func{g}{\Omega\subseteq \R^3}{\R}$ un campo escalar, ambos diferenciables, se tiene que
    \[
        \iiint_{\Omega} (F \cdot \nabla g + g \dive(F))\, dV
        = \varoiint_{\partial \Omega}  g F \cdot dS.
    \]
\end{prop}

\begin{prop}[Fórmula de Green]
    Sea $\Omega$ una región acotada de $\R^3$ tal que $\partial \Omega$ es una superficie suave, cerrada y parametrizada con orientación hacia fuera de $\Omega$. Sean $\func{f,g}{\Omega}{\R}$ campos escalares dos veces diferenciables, se tiene que
    \[
        \iiint_\Omega (\nabla f\cdot \nabla g + f \Delta g) \, dV
        =
        \varoiint_{\partial \Omega} (f\nabla g) \cdot dS.
    \]
\end{prop}

\begin{prop}
    Sea $\Omega$ una región acotada de $\R^3$ tal que $\partial \Omega$ es una superficie suave, cerrada y parametrizada con orientación hacia fuera de $\Omega$. Sean $\func{F,G}{\Omega}{\R^3}$ campos vectoriales dos veces diferenciables, se tiene que
    \[
        \iiint_{\Omega} (G \cdot \rot(F) - F \cdot \rot(G))\, dV
        = \varoiint_{\partial \Omega}  (F\times G) \cdot dS.
    \]
\end{prop}


\end{document}
